\subsection{Autofunzioni degli operatori}
So che \( \hat{q} \ket{q} = q \ket{q} \), dunque
\[
\braket{q^\prime|q} = \psi_q(q^\prime) = \delta(q - q^\prime) \; \rightarrow \; \text{autofunzione della posizione.}
\]

Per l'impulso voglio che esista \( \ket{k} \) tale che \( \hat{p} \ket{k} = \hbar k \ket{k} \). Scelgo la base dei \( \ket{q} \) e faccio
\[
\braket{q|\hat{p}|k} = \hbar k \braket{q|k}, \quad \text{con } \braket{q|k} = \psi_k(q) \quad \Rightarrow
\]
\[
\Rightarrow \quad -i \hbar \frac{d}{dq} \psi_k(q) = \hbar k \psi_k(q) \quad \Rightarrow \quad \psi_k(q) = N_k e^{ikq}.
\]
Poiché
\[
\int_{-\infty}^{+\infty} \left| \psi_k(q) \right|^2 \, dq = |N_k|^2 \int_{-\infty}^{+\infty} 1 \, dq,
\]
non posso normalizzare in senso proprio. Coerentemente con le autofunzioni della posizione, ho \( \braket{k|k^\prime} = \delta(k - k^\prime) \), cioè
\begin{align*}
\braket{k|k^\prime} &= \int_{-\infty}^{+\infty} \braket{k^\prime|q} \braket{q|k} \, dq = |N_k|^2 \int_{-\infty}^{+\infty} e^{-i k^\prime q} e^{ikq} \, dq = \\
&= |N_k|^2 \int_{-\infty}^{+\infty} e^{i (k - k^\prime) q} \, dq = \delta(k^\prime - k).
\end{align*}
Allora
\begin{align*}
1 &= \int_{-\infty}^{+\infty} \braket{k|k'} \, dk = \int_{-\infty}^{+\infty} \delta(k - k') \, dk = \\
&= \int_{\mathbb{R}^2} |N_k|^2 e^{i (k - k^\prime) q} \, dq \, dk = \\
&= \lim_{\lambda \rightarrow +\infty} \int_{-\infty}^{+\infty} |N_k|^2 \left[ \frac{e^{i (k - k^\prime) q}}{i (k - k^\prime)} \right]_{-\lambda}^{+\lambda} \, dk = \\
&= \lim_{\lambda \rightarrow +\infty} \int_{-\infty}^{+\infty} |N_k|^2 \frac{e^{i (k - k^\prime) \lambda} - e^{-i (k - k^\prime) \lambda}}{i (k - k^\prime)} \, dk = \\
&= 2 \pi |N_k|^2 \qquad \Rightarrow \qquad |N_k|^2 = \frac{1}{2\pi}
\end{align*}
\[
\psi_k(q) = \frac{1}{\sqrt{2\pi}} e^{ikq}
\]
e ridefinisco
\[
\delta(x) = \frac{1}{2\pi} \int_{-\infty}^{+\infty} e^{ipx} \, dp.
\]
Dalla definizione di aggiunto,
\[
\braket{q|k} = \frac{1}{\sqrt{2\pi}} e^{ikq}, \qquad \braket{k|q} = \frac{1}{\sqrt{2\pi}} e^{-ikq}
\]
quindi analogo ad \( U^\dag U = \mathbbm{1} \), cioè matrici di cambio di base
\[
\braket{q|\psi} = \psi(q), \qquad \braket{k|\psi} = \tilde{\psi}(k)
\]
legate da
\begin{align*}
\braket{k|\psi} &= \tilde{\psi}(k) = \int_{-\infty}^{+\infty} \braket{k|q} \braket{q|\psi} \, dq = \\
&= \int_{-\infty}^{+\infty} \frac{e^{-ikq}}{\sqrt{2\pi}} \psi(q) \, dq \\
\braket{q|\psi} &= \int_{-\infty}^{+\infty} \braket{q|k} \braket{k|\psi} \, dk = \int_{-\infty}^{+\infty} \frac{e^{ikq}}{\sqrt{2\pi}} \tilde{\psi}(k) \, dk.
\end{align*}
N.B.: in generale, \( \tilde{\psi} \neq \psi \).

Poiché
\[
\psi(q) = \braket{q|\psi} = \int_{-\infty}^{+\infty} \braket{q|k} \braket{k|\psi} \, dk = \frac{1}{\sqrt{2\pi}} \int_{-\infty}^{+\infty} e^{-ikq} \psi(k) \, dk
\]
\[
\psi(k) = \braket{k|\psi} = \int_{-\infty}^{+\infty} \braket{k|q} \braket{q|\psi} \, dq = \frac{1}{\sqrt{2\pi}} \int_{-\infty}^{+\infty} e^{ikq} \psi(q) \, dq
\]
ottengo che \( \ket{q} \) e \( \ket{k} \) sono \textbf{linearmente dipendenti}.

Gli elementi di matrice
\[
\braket{k|\hat{p}|\psi} = \hbar k \psi(k).
\]
N.B.: \( \hat{p}\ket{p} = p \ket{p} \Leftrightarrow \braket{p|\hat{p}|\psi}  = p \psi(p) \) con \( \psi(p) = \braket{p|\psi} \), ma allora
\[
\braket{q|p} = \frac{1}{\sqrt{2\pi}} e^{i \frac{pq}{\hbar}}.
\]

\begin{align*}
\braket{k|\hat{q}|\psi} &= \int_{-\infty}^{+\infty} \braket{k|\hat{q}|q} \braket{q|\psi} \, dq = \\
&= \int_{-\infty}^{+\infty} q \braket{k|q} \psi(q) \, dq = \\
&= \int_{-\infty}^{+\infty} q \frac{e^{-ikq}}{\sqrt{2\pi}} \psi(q) \, dq = \frac{1}{\sqrt{2\pi}} \int_{-\infty}^{+\infty} i \frac{d}{dk} e^{-ikq} \psi(q) \, dq = \\
&= i \frac{1}{\sqrt{2\pi}} \frac{d}{dk} \underbrace{\int_{-\infty}^{+\infty} e^{-ikq} \psi(q) \, dq}_{\psi(k)} = i \frac{1}{\sqrt{2\pi}}\frac{d}{dk} \psi(k).
\end{align*}